\documentclass[11pt, a4paper]{article}
\usepackage[margin=1in]{geometry}
\usepackage[T1]{fontenc}
\usepackage{lmodern}
\usepackage{microtype}
\usepackage{booktabs}
\usepackage{tabularx}
\usepackage{longtable}
\usepackage{amsmath, amssymb}
\usepackage{enumitem}
\usepackage{xcolor}
\usepackage{hyperref}
\usepackage{fancyvrb}
\usepackage{float}
\usepackage{caption}

\definecolor{accentblue}{HTML}{2563EB}
\definecolor{muted}{HTML}{64748B}
\definecolor{codebg}{HTML}{F8FAFC}

\hypersetup{
  colorlinks=true,
  linkcolor=accentblue,
  urlcolor=accentblue,
  citecolor=accentblue,
}

\setlength{\parindent}{0pt}
\setlength{\parskip}{6pt}

% Section formatting
\usepackage{titlesec}
\titleformat{\section}{\large\bfseries\sffamily}{\thesection.}{0.5em}{}[\vspace{-4pt}\rule{\textwidth}{0.4pt}]
\titleformat{\subsection}{\normalsize\bfseries\sffamily}{\thesubsection}{0.5em}{}
\titleformat{\subsubsection}{\normalsize\itshape}{\thesubsubsection}{0.5em}{}

\title{%
  \sffamily\bfseries
  Perception Agent Benchmark:\\[4pt]
  \Large A Representation Ablation Study for\\
  Vision-Based Embryo Stage Classification
}
\author{%
  Gently Project, Shroff Lab, Janelia Research Campus
}
\date{February 2026}

\begin{document}
\maketitle

\begin{abstract}
The Gently microscopy platform uses Claude's vision capabilities to classify
\textit{C.\ elegans} embryo developmental stages from light-sheet 3D volumes in real
time. The fundamental challenge is representational: Claude cannot see in 3D, so
volumetric data must be projected into 2D images. The current approach (three
orthogonal max-intensity projections) achieves 33\% exact accuracy and 69\%
adjacent accuracy. Even the human annotator noted difficulty classifying from
projections alone. This document defines a systematic benchmark to determine
whether the representation is the bottleneck, and proposes controlled ablation
experiments across four representation strategies, reference example counts, and
tool-use configurations. The dataset (3D volumes, ground truth annotations,
baseline results) is publicly available at
\url{https://huggingface.co/datasets/pskeshu/gently-perception-benchmark}.
\end{abstract}

\tableofcontents
\vspace{12pt}

% ═══════════════════════════════════════════════════════════════
\section{The Problem}
% ═══════════════════════════════════════════════════════════════

\subsection{Biological context}

\textit{C.\ elegans} embryo development proceeds through a continuous
morphological transformation: a fertilized egg undergoes gastrulation, elongation,
folding, and coiling until a fully formed larva hatches from the eggshell. This
process takes approximately 13--14 hours at 20\textdegree C and is conventionally
described in discrete stages:

\begin{center}
\small
\begin{tabular}{@{}lll@{}}
\toprule
\textbf{Stage} & \textbf{Key feature} & \textbf{Duration (approx.)} \\
\midrule
early    & Symmetric oval, dividing cells         & $\sim$350 min \\
bean     & End asymmetry, central constriction     & $\sim$20 min \\
comma    & Ventral indentation, C-shape            & $\sim$30 min \\
1.5-fold & Body folding back on itself             & $\sim$30 min \\
2-fold   & Body doubled back completely            & $\sim$45 min \\
pretzel  & Tight coil, 3+ body segments            & $\sim$300 min \\
hatching & Active emergence from shell             & minutes \\
hatched  & Free L1 larva                           & --- \\
\bottomrule
\end{tabular}
\end{center}

These stages are \emph{bins imposed on a continuum}. The transitions are gradual,
not instantaneous. Any classification system must contend with this inherent
ambiguity.

\subsection{Imaging modality}

We capture the embryo using a dual-view selective plane illumination microscope
(diSPIM). The light-sheet creates a 2D illumination plane that is scanned through
the depth of the specimen, producing a 3D volume (z-stack). The fluorescence
signal comes from a histone marker labeling all nuclei.

The raw data is inherently \textbf{3D}: a volume of shape $(Z, Y, X)$ where $Z$
is typically 40--80 slices and $Y, X$ are the lateral dimensions.

\subsection{The fundamental gap}

Claude is a 2D vision model. It cannot consume a 3D volume directly. The entire
perception system therefore hinges on \emph{how we project the 3D volume into 2D
images} that the model can see.

This creates a representational bottleneck with several failure modes:

\begin{enumerate}[leftmargin=*, itemsep=2pt]
  \item \textbf{Projections destroy depth structure.} Max-intensity projections
    collapse an entire axis into a single plane. Folding and coiling, which are
    3D spatial relationships, merge into ambiguous 2D patterns when overlapping
    structures are superimposed.

  \item \textbf{The projection choice is arbitrary.} Why XY, YZ, and XZ? Why
    max-intensity rather than mean, sum, or alpha compositing? Each choice
    destroys different information. The system prompt compensates by instructing
    the model to ``ALWAYS ANALYZE XZ VIEW,'' an admission that the
    representation is fragile.

  \item \textbf{The model is not a microscopist.} Claude was trained on internet
    images, not fluorescence microscopy. The perception prompt bridges this gap
    with text descriptions (``peanut shape'', ``horizontal figure-8''), asking
    the model to map novel visual patterns to verbal descriptions. This mapping
    is unvalidated.

  \item \textbf{Few-shot examples are thin.} 2 reference images per stage
    $\times$ 7 stages = 14 total images to represent the full morphological
    diversity of embryo development. Every embryo looks different; orientation
    varies; imaging conditions vary.

  \item \textbf{Verification doesn't add information.} When Phase~1 is uncertain,
    spawning Haiku subagents to examine the \emph{same projections} averages over
    the same noise rather than resolving the underlying ambiguity.
\end{enumerate}

The annotator's note is telling:
\begin{quote}
\itshape
``I did find it difficult just using the max projections alone, as sometimes the
embryo would orient itself in a way that made it hard to determine the location of
the nose or tail.'' (Ryan, biologist, December 2024)
\end{quote}

If a domain expert struggles with the representation, a general-purpose vision model will
struggle more.

% ═══════════════════════════════════════════════════════════════
\section{Current Perception Architecture}
% ═══════════════════════════════════════════════════════════════

\subsection{System overview}

The perception agent is a two-phase pipeline:

\begin{Verbatim}[fontsize=\small, frame=single, framesep=8pt, baselinestretch=0.95]
  Caller (Timelapse Orchestrator / Benchmark Runner)
    |
    |  embryo_id, timepoint, image_b64, volume?
    v
  PerceptionManager            (one session per embryo)
    |
    |  engine.perceive(image, session, timepoint, volume?)
    v
  PerceptionEngine             MODEL: claude-opus-4-6
    |
    |  _build_prompt():
    |    STATIC (cached 1h):  system prompt + 14 reference images
    |    DYNAMIC (per-call):  last 3 observations, temporal context,
    |                         current 3-view image
    |
    |  _run_reasoning_loop()   (max 5 tool-use iterations)
    |    |
    |    +-- end_turn --> parse JSON --> calibrate confidence
    |    |                               --> PerceptionResult
    |    +-- tool_use:
    |          +-- view_previous_timepoint  (cached image)
    |          +-- view_embryo              (3D alpha-composite)
    |          +-- request_verification --> Phase 2
    |
    v  (if verification triggered, confidence < 0.7)
  VerificationEngine           SUBAGENT MODEL: claude-haiku-4-5-20251001
    |
    |  asyncio.gather: up to 3 parallel subagents
    |    Subagent A: "stage_a vs stage_b"  (focused comparison)
    |    Subagent B: "stage_c vs stage_d"
    |    Subagent C: "stage_e vs stage_f"
    |
    |  _aggregate_results(): confidence-weighted voting
    v
  PerceptionResult
    stage, confidence, is_hatching, is_transitional,
    observed_features, contrastive_reasoning,
    reasoning_trace, verification_result
\end{Verbatim}

\noindent\small\textit{The \texttt{?} on \texttt{volume} indicates
\texttt{Optional[np.ndarray]} in the API; the raw 3D volume is passed when
available (live imaging), but the system can also operate from pre-rendered 2D
images alone (offline/benchmark mode).}
\normalsize

\subsection{Current performance}

The most recent benchmark run (Sonnet 4.5, December 2024) on session
\texttt{59799c78} with 4 embryos and 737 total timepoints:

\begin{center}
\begin{tabular}{@{}lr@{}}
\toprule
\textbf{Metric} & \textbf{Value} \\
\midrule
Exact accuracy         & 33.2\% \\
Adjacent accuracy      & 69.5\% \\
Total predictions      & 737 \\
Model                  & \texttt{claude-sonnet-4-5-20250929} \\
Tools                  & all enabled \\
Verification           & enabled \\
Reference examples     & 2 per stage \\
\bottomrule
\end{tabular}
\end{center}

33\% exact accuracy means the model gets the correct stage only one-third of the
time. 69\% adjacent accuracy means it is within one stage about two-thirds of the
time. This is insufficient for reliable autonomous microscopy.

\subsection{Baseline analysis}

Breaking down the baseline results reveals systematic failure patterns.

\subsubsection{Per-stage accuracy}

\begin{center}
\small
\begin{tabular}{@{}lrrr@{}}
\toprule
\textbf{Stage} & \textbf{Samples} & \textbf{Accuracy} & \textbf{Most common error} \\
\midrule
early    & 157 & 40.8\% & predicted as bean (93/157) \\
bean     &  24 & 79.2\% & predicted as comma (5/24) \\
comma    &  27 & 44.4\% & predicted as bean (15/27) \\
1.5-fold &  49 & 28.6\% & predicted as comma (33/49) \\
2-fold   &  79 & 67.1\% & predicted as comma (16/79) \\
pretzel  & 401 & 20.7\% & predicted as 2-fold (110/401) \\
\bottomrule
\end{tabular}
\end{center}

Two patterns dominate: the model is systematically late on early$\to$bean (calls
59\% of early timepoints ``bean''), and systematically early on pretzel (calls
27\% of pretzel timepoints ``2-fold''). The 1.5-fold stage is almost entirely
missed, with 67\% of predictions falling on comma instead.

\subsubsection{Confusion matrix}

\begin{center}
\small
\begin{tabular}{@{}l*{6}{r}@{}}
\toprule
& \multicolumn{6}{c}{\textbf{Predicted}} \\
\cmidrule(l){2-7}
\textbf{True} & early & bean & comma & 1.5f & 2f & pretzel \\
\midrule
early    & \textbf{64} & 93  &  0 &  0 &   0 &  0 \\
bean     &  0 & \textbf{19} &  5 &  0 &   0 &  0 \\
comma    &  0 & 15 & \textbf{12} &  0 &   0 &  0 \\
1.5-fold &  0 &  1 & 33 & \textbf{14} &   1 &  0 \\
2-fold   &  0 &  0 & 16 & 10 &  \textbf{53} &  0 \\
pretzel  &  0 &  0 &  0 &  0 & 110 & \textbf{83} \\
\bottomrule
\end{tabular}
\end{center}

The matrix shows a strong upper-diagonal bias: errors almost always predict an
\emph{earlier} stage than the true one, suggesting the model is slow to recognize
transitions.

\subsubsection{Calibration}

The model's confidence is poorly calibrated. Mean confidence when correct (0.867)
is nearly identical to mean confidence when wrong (0.857), yielding an Expected
Calibration Error (ECE) of 0.524. The model does not know when it is wrong.

\subsubsection{Tool use}

Tools were invoked on 16.6\% of predictions (122/737 total tool calls).
Accuracy \emph{with} tool use was 23.0\%, compared to 35.1\% \emph{without}.
This suggests that tool invocation correlates with difficult cases and does not
resolve the underlying ambiguity, consistent with the hypothesis that
verification on the same representation cannot add new information.

\subsection{The representation pipeline (current)}

The volume-to-image conversion is implemented in
\texttt{projection.py} and \texttt{testset.py}:

\begin{enumerate}[leftmargin=*, itemsep=2pt]
  \item Load 3D volume from TIFF: shape $(Z, Y, X)$
  \item Auto-crop using center-of-mass (95th percentile threshold)
  \item Generate three max-intensity projections:
    \begin{itemize}[itemsep=0pt]
      \item XY: \texttt{np.max(volume, axis=0)}, looking down
      \item YZ: \texttt{np.max(volume, axis=2)}, looking from side
      \item XZ: \texttt{np.max(volume, axis=1)}, looking from front
    \end{itemize}
  \item Composite into a single image: $[\text{XY}|\text{YZ}]$ top row,
    $[\text{XZ}]$ bottom row
  \item Normalize (1st--99th percentile), resize to $\leq$1500px, encode as
    base64 JPEG
\end{enumerate}

This is the representation we aim to challenge in the benchmark.

% ═══════════════════════════════════════════════════════════════
\section{Benchmark Design}
% ═══════════════════════════════════════════════════════════════

\subsection{Central question}

\textbf{Is the 2D representation the bottleneck, or is the model the bottleneck?}

If different representations of the same 3D data produce significantly different
accuracy, the representation is the bottleneck. If accuracy is uniformly low
regardless of representation, the problem is deeper (model capability, few-shot
inadequacy, or task formulation).

\subsection{Experimental structure}

Three sequential experiments, each informed by the results of the previous:

\begin{center}
\small
\begin{tabularx}{\textwidth}{@{}clXl@{}}
\toprule
\textbf{Exp} & \textbf{Question} & \textbf{Independent variable}
  & \textbf{Conditions} \\
\midrule
1 & Is representation the bottleneck?
  & Volume $\to$ image strategy
  & 4 \\
2 & Do few-shot examples help?
  & Number of reference images per stage
  & 3 \\
3 & Do tools/verification help?
  & Tool access configuration
  & 3 \\
\bottomrule
\end{tabularx}
\end{center}

\textbf{Held constant} across all conditions within an experiment: model
(configurable, Opus or Sonnet), ground truth, embryo set, timepoint set.

\subsection{Experiment 1: Representation ablation}

Four representation strategies, each taking the same 3D volume as input:

\subsubsection{A. Three-View MIP (baseline)}

The current approach. Three orthogonal max-intensity projections composited into a
single image. The model sees one image containing XY, YZ, and XZ views.

\textit{Strengths:} Compact (single image), computationally cheap.\\
\textit{Weaknesses:} Collapses depth, merges overlapping structures.

\subsubsection{B. Z-Slice Montage}

Show the model a grid of individual z-slices (e.g., 16 evenly-spaced slices in a
$4 \times 4$ grid). Each tile is labeled with its z-index. The model can see
which structures appear at which depths; stacking vs.\ side-by-side arrangement
becomes directly visible.

\textit{Strengths:} Preserves depth structure, no information collapse, closest
to what a microscopist sees when scrolling.\\
\textit{Weaknesses:} Lower per-tile resolution, may be harder for the model to
integrate mentally.

\subsubsection{C. Multi-Angle Rendered Views}

Render the volume from 6 canonical viewing angles using alpha compositing
(reusing the existing \texttt{render\_volume\_view()} function):

\begin{center}
\small
\begin{tabular}{@{}lrrl@{}}
\toprule
\textbf{View} & $r_x$ & $r_y$ & \textbf{Perspective} \\
\midrule
Top-down   & 0\textdegree   & 0\textdegree    & Default \\
Tilted     & 45\textdegree  & 0\textdegree    & Forward tilt \\
Right side & 0\textdegree   & 90\textdegree   & Side view \\
Left side  & 0\textdegree   & $-$90\textdegree & Opposite side \\
Oblique R  & 45\textdegree  & 45\textdegree   & Angled right \\
Front      & 90\textdegree  & 0\textdegree    & Front-on \\
\bottomrule
\end{tabular}
\end{center}

Views composited into a $2 \times 3$ grid.

\textit{Strengths:} Shows 3D shape from multiple perspectives, depth-aware
rendering.\\
\textit{Weaknesses:} Each view is still a 2D projection (albeit with depth
cues), computationally heavier.

\subsubsection{D. Classical Features + Single MIP}

Compute quantitative 3D morphological features from the segmented volume and pass
them as structured text alongside a single top-down projection:

\begin{itemize}[itemsep=1pt]
  \item Bounding box aspect ratios ($X/Y$, $X/Z$, $Y/Z$)
  \item Elongation (PCA major/minor axis ratio)
  \item Compactness ($\text{volume} / \text{bounding\_box\_volume}$)
  \item Sphericity ($36\pi V^2 / S^3$)
  \item Number of connected components per z-level
  \item Centroid position relative to bounding box
  \item Moments of inertia ratios
\end{itemize}

\textit{Strengths:} Offloads 3D understanding to computation; the model reasons over
pre-computed features rather than extracting them from unfamiliar images.\\
\textit{Weaknesses:} Depends on segmentation quality; loses visual nuance the model
might catch.

\subsection{Experiment 2: Reference example count}

Using the best representation from Experiment~1:

\begin{center}
\small
\begin{tabular}{@{}lll@{}}
\toprule
\textbf{Condition} & \textbf{Examples/stage} & \textbf{What it tests} \\
\midrule
Zero-shot  & 0 & Can the model classify from the prompt alone? \\
One-shot   & 1 & Minimal visual reference \\
Two-shot   & 2 & Current default \\
\bottomrule
\end{tabular}
\end{center}

The reference images consume $\sim$21{,}000 tokens (cached). If zero-shot
performs comparably, these tokens are wasted. If more examples help significantly,
the example store needs expansion.

\subsection{Experiment 3: Tool and verification ablation}

Using the best representation and example count:

\begin{center}
\small
\begin{tabular}{@{}llll@{}}
\toprule
\textbf{Condition} & \textbf{Tools} & \textbf{Verification} & \textbf{What it tests} \\
\midrule
Full       & on  & on  & Current system (max capability) \\
Tools only & on  & off & Does interactive exploration help? \\
Minimal    & off & off & Raw single-pass classification \\
\bottomrule
\end{tabular}
\end{center}

% ═══════════════════════════════════════════════════════════════
\section{API Cost Estimates}
% ═══════════════════════════════════════════════════════════════

\subsection{Token budget per call}

\begin{center}
\small
\begin{tabular}{@{}lrl@{}}
\toprule
\textbf{Component} & \textbf{Tokens} & \textbf{Caching} \\
\midrule
System prompt (text)                & $\sim$800     & cached (1h TTL) \\
Reference images ($14 \times 1500$) & $\sim$21{,}000 & cached (1h TTL) \\
Previous observations               & $\sim$200     & uncached \\
Temporal context                     & $\sim$200     & uncached \\
Current image                        & $\sim$1{,}500 & uncached \\
\midrule
\textbf{Total input}                & $\sim$23{,}700 & 22K cached + 1.9K fresh \\
\textbf{Output} (reasoning + JSON)  & $\sim$1{,}000  & --- \\
\bottomrule
\end{tabular}
\end{center}

Representation-specific image token variations: slice montage $\sim$3{,}000;
multi-angle $\sim$4{,}000; features+MIP $\sim$1{,}300 (smaller image + 300 tokens
text).

\subsection{Per-call and per-condition costs}

\begin{center}
\small
\begin{tabular}{@{}lrrrr@{}}
\toprule
\textbf{Model} & \textbf{Cache read} & \textbf{Fresh input}
  & \textbf{Output} & \textbf{Per call} \\
\midrule
Opus 4.6
  & 22K $\times$ \$0.50/M & 2K $\times$ \$5/M
  & 1K $\times$ \$25/M & \textbf{\$0.046} \\
Sonnet 4.6
  & 22K $\times$ \$0.30/M & 2K $\times$ \$3/M
  & 1K $\times$ \$15/M & \textbf{\$0.028} \\
Haiku 4.5 (verif.)
  & --- & 5K $\times$ \$1/M
  & 0.5K $\times$ \$5/M & \textbf{\$0.008} \\
\bottomrule
\end{tabular}
\end{center}

\subsection{Total experiment costs}

Test data: 4 embryos, 737 timepoints total.

\begin{center}
\small
\begin{tabular}{@{}lcrr@{}}
\toprule
\textbf{Experiment} & \textbf{Conditions}
  & \textbf{Opus 4.6} & \textbf{Sonnet 4.6} \\
\midrule
1: Representation    & 4 & \$136 & \$83 \\
2: Example count     & 3 & \$102 & \$62 \\
3: Tools/verification & 3 & \$102 & \$62 \\
\midrule
\textbf{Total}       & \textbf{10} & \textbf{\$339} & \textbf{\$207} \\
\bottomrule
\end{tabular}
\end{center}

\textbf{Quick validation run} (30 timepoints/embryo $\approx$ 120 total):
Opus \$55, Sonnet \$34.

\textbf{Batch API option:} The Batch API provides a 50\% discount for
non-real-time workloads, reducing the full experiment to $\sim$\$170 (Opus) /
$\sim$\$104 (Sonnet).

Verification subagent overhead (Haiku 4.5, $\sim$15\% of calls triggered): adds
$\sim$5\% to totals.

% ═══════════════════════════════════════════════════════════════
\section{Metrics and Analysis}
% ═══════════════════════════════════════════════════════════════

\subsection{Primary metrics (per condition)}

\begin{enumerate}[leftmargin=*, itemsep=2pt]
  \item \textbf{Exact accuracy:} fraction of predictions matching ground truth
  \item \textbf{Adjacent accuracy:} fraction within $\pm 1$ stage
  \item \textbf{Per-stage accuracy:} accuracy broken down by ground truth stage
  \item \textbf{Confusion matrix:} predicted vs.\ ground truth counts
  \item \textbf{Expected Calibration Error (ECE):} do confidence scores reflect
    true accuracy?
  \item \textbf{Backward transition rate:} how often does the predicted stage
    regress (biologically impossible)?
\end{enumerate}

\subsection{Comparison metrics (across conditions)}

\begin{enumerate}[leftmargin=*, itemsep=2pt]
  \item \textbf{$\Delta$accuracy by representation:} which representation wins,
    by how much?
  \item \textbf{Stage-specific $\Delta$:} which stages benefit most from
    better representations?
  \item \textbf{Accuracy vs.\ cost:} is the cheapest representation good enough?
  \item \textbf{Tool use impact:} accuracy with vs.\ without tools, stratified
    by stage
\end{enumerate}

\subsection{Planned figures}

The benchmark report will include the following visualizations:

\begin{center}
\small
\begin{tabularx}{\textwidth}{@{}clX@{}}
\toprule
\textbf{Fig.} & \textbf{Type} & \textbf{Content} \\
\midrule
1 & Grouped bar chart
  & Exact \& adjacent accuracy by representation (Exp 1) \\
2 & Heatmaps ($4\times$)
  & Confusion matrix for each representation \\
3 & Line plot
  & Per-stage accuracy across representations (one line per repr) \\
4 & Grouped bar chart
  & Accuracy by example count: 0-shot vs 1-shot vs 2-shot (Exp 2) \\
5 & Stacked bar chart
  & Accuracy with/without tools and verification (Exp 3) \\
6 & Scatter plot
  & Confidence vs.\ actual accuracy (calibration) per representation \\
7 & Timeline plot
  & Predicted stage vs.\ ground truth over time for one embryo,
    per representation (qualitative comparison) \\
8 & Cost/accuracy Pareto
  & Each condition plotted as (cost, accuracy) to identify the
    efficient frontier \\
\bottomrule
\end{tabularx}
\end{center}

% ═══════════════════════════════════════════════════════════════
\section{Validation}
% ═══════════════════════════════════════════════════════════════

\begin{enumerate}[leftmargin=*, itemsep=4pt]
  \item \textbf{Unit:} Each representation strategy produces a valid base64 JPEG
    from a test volume.
  \item \textbf{Smoke:} Run benchmark with \texttt{-{}-max-timepoints 3
    -{}-embryo embryo\_1} for each representation; verify completion and valid
    JSON output.
  \item \textbf{Quick ablation:} All 4 representations on 30
    timepoints/embryo with Sonnet ($\sim$\$5 total); verify comparison report
    generates correctly.
  \item \textbf{Full Experiment 1:} All 4 representations on all 737 timepoints.
    Analyze results before proceeding to Experiments 2 and 3.
\end{enumerate}

% ═══════════════════════════════════════════════════════════════
\section{Expected Outcomes}
% ═══════════════════════════════════════════════════════════════

\subsection{If representation is the bottleneck}

We expect the z-slice montage or multi-angle views to outperform the baseline,
particularly on fold stages (1.5-fold, 2-fold, pretzel) where 3D spatial
relationships are critical. The early-stage transitions (early/bean/comma) may
show less improvement since they rely more on 2D shape cues.

\emph{Next step:} Invest in better representations. Explore hybrid approaches
(e.g., combining montage with rendered views) and optimized prompt strategies
for the best-performing representation.

\subsection{If representation is not the bottleneck}

If all representations perform similarly poorly, the problem is at the model or
task level:
\begin{itemize}[itemsep=2pt]
  \item The few-shot examples may be inadequate (Experiment~2 will test this)
  \item The model may lack the visual priors for fluorescence microscopy
  \item The stage boundaries may be too ambiguous for vision-based classification
\end{itemize}

\emph{Next step:} Evaluate whether a more capable model (Opus vs.\ Sonnet)
changes the picture, and whether prompt or example improvements can close the
gap.

\end{document}
